\documentclass[../../../include/open-logic-section]{subfiles}

\begin{document}

\olfileid{sth}{story}{urelements}
\olsection{मूलघटक असावेत का?}

पुढील काही प्रकरणांत आपण संचांच्या संचयी-क्रमिक संकल्पनेतून स्वयंसिद्धके मिळवण्याचा प्रयत्न करू. पण पुढे जाण्याआधी \emph{मूलघटकांबद्दल} आणखी काही सांगणे गरजेचे आहे.

\olref[approach]{sec} मधील चित्रात आपण जुन्या \emph{संचांपासून} नवे संचच तयार करत होतो. पण काही \emph{संच नसलेल्या} वस्तू---गायी, डुकरे, वाळूचे कण किंवा इतर काही---संचांचे !!{element}s असू शकतात, असे मानायचे असेल. तसे असल्यास सुरुवातीला काही मूलभूत वस्तू, म्हणजे \emph{मूलघटक}, उपलब्ध मानू. मग प्रत्येक टप्प्या~$S$ वर त्या मूलघटकांपासून किंवा $S$ पूर्वीच्या टप्प्यांवर तयार झालेल्या संचांपासून बनणारे ``सर्व शक्य'' संच तयार होतात, असे म्हणू; यात दोन्ही प्रकारच्या वस्तू एकत्रही असू शकतात. अशा वेळी चित्र साधारण असे दिसेल:
\begin{center}
	\begin{tikzpicture}[scale=0.6]
	\tikzset{cut_here/.style={densely dotted}}
	\draw (-4,6) -- (-1, 0) -- (1,0) -- (4,6);
	\draw[cut_here] (-5,8)--(-4,6);
	\draw[cut_here] (5,8)--(4,6);
	\draw(-1.5, 1)--(1.5, 1);
	\draw(-2, 2)--(2, 2);
	\draw(-2.5, 3)--(2.5,3);
	\draw(-3, 4)--(3, 4);
	\draw(-3.5, 5)--(3.5, 5);
	\draw(-4, 6)--(4, 6);
	\node[label] at (2, 0) {\small 0};
	\node[label] at (2.5, 1) {\small 1};
	\node[label] at (3, 2) {\small 2};
	\node[label] at (3.5, 3) {\small 3};
	\node[label] at (4, 4) {\small 4};
	\node[label] at (4.5, 5) {\small 5};
	\node[label] at (5, 6) {\small 6};
	\end{tikzpicture}
\end{center}
आता आपल्याला एक निर्णय घ्यायचा आहे: \emph{मूलघटकांना परवानगी द्यावी का?}

तात्त्विक दृष्टीने आपल्या सिद्धान्तात मूलघटक समाविष्ट करणे सयुक्तिक आहे. त्यामागचे मुख्य कारण म्हणजे संच उपपत्ती \emph{लागू करता येण्याजोगी} बनवणे. हे स्पष्ट करण्यासाठी \olref[sfr][siz][]{chap} आठवा: दोन संच $A$ आणि~$B$ यांचा आकार समान, म्हणजे $\cardeq{A}{B}$, तेव्हा आणि केवळ तेव्हाच त्यांच्यात एकास-एक व आच्छादक फलन असते. आता शेतातल्या गायी आणि गोठ्यातली डुकरे हे दोन्ही संच बनवत असतील, तर ``गायी जितक्या आहेत तितकीच डुकरे आहेत'' या विधानाचा संच उपपत्तीच्या आधारे विचार करता येतो. पण मूलघटकांना मनाई केली आणि गायी व डुकरे यांचे संच \emph{बनतच नाहीत} असे मानले, तर हा विचार करता येणार नाही. खरे तर संचांशिवाय इतर कशालाही संच उपपत्ती लागू करण्याचा सरळ मार्ग आपल्याकडे उरणार नाही. (मूलघटक समाविष्ट करण्याची आणखी कारणे \citealt[pp.~vi, 24,
50--1]{Potter2004} येथे पाहा.)

मात्र गणितात मूलघटकांना परवानगी देणे फारसे प्रचलित नाही. त्याचे एक कारण म्हणजे मूलघटकांशिवाय संच उपपत्ती मांडणे \emph{अगदी किंचित} सोपे असते. पण कधीकधी त्यांना वगळण्याचे अधिक लक्षवेधी समर्थनही आढळते:
\begin{quote}
  संच उपपत्ती हा गणिताचा पाया आहे या मतानुसार, केवळ संचांविषयी बोलून आपल्याला संपूर्ण गणित व्यक्त करता आले पाहिजे. म्हणून आपल्या चलांची व्याप्ती गायी आणि डुकरे यांसारख्या वस्तूंवर असता कामा नये.
  %But if $C$ is a cow, $\{C\}$ is a set, but not a legitimate mathematical object.
  \citep[p.~8]{Kunen1980}
\end{quote}
म्हणून उपयोज्यतेला महत्त्व दिल्यास मूलघटक \emph{समाविष्ट} करावेसे वाटतील; पण गणित शुद्ध संच उपपत्तीत आणण्याच्या पायाभूत ध्येयाला महत्त्व दिल्यास त्यांना \emph{वगळावेसे} वाटेल. थोडा आळसही मूलघटक वगळण्याकडे झुकवतो.

आपण सर्वांत सोपा मार्ग निवडू. पण त्यामागे काही प्रमाणात शैक्षणिक कारणही आहे. संच उपपत्तीच्या तत्त्वज्ञानाचा अभ्यास सुरू करण्यासाठी तुम्हाला तिच्या ज्या मूलभूत गोष्टी लागतील, त्यांची ओळख करून देणे हे आपले उद्दिष्ट आहे. या विषयावरील बहुतेक तात्त्विक साहित्य मूलघटक \emph{नसलेल्या} संच उपपत्तीची चर्चा करते. म्हणून त्या साहित्याशी साधारण जुळून राहिल्यास हे पुस्तक अधिक उपयोगी ठरेल.

\end{document}
